% !TEX root = ../thesis.tex

\chapter{Koncept riešenia a odhad veľkosti systému}\label{ch:concept-and-size}

Táto kapitola prepája doteraz definované časti práce do jedného realizovateľného konceptu. Cieľom je zjednotiť víziu riešenia, určiť hlavné moduly systému, odhadnúť ich náročnosť a pripraviť podklad pre prvú iteráciu návrhu a implementácie.

\section{Koncept a vízia riešenia}

Navrhované riešenie predstavuje webový informačný systém pre správu distribúcie pitnej a úžitkovej vody, ktorý centralizuje evidenciu zariadení, revízií, servisných zásahov a súvisiacich dokumentov. Víziou je odstrániť fragmentáciu údajov medzi oddeleniami a nahradiť manuálne evidencie jednotným, auditovateľným a prehľadným pracovným prostredím.

Z pohľadu používateľa je systém navrhnutý ako roly orientovaná platforma:

\begin{itemize}
\item správca siete plánuje revízie, sleduje termíny a stav zariadení,
\item technik vykonáva servisné zásahy a zaznamenáva priebeh opráv,
\item revízny pracovník dokumentuje výsledky kontrol a nápravné opatrenia,
\item manažér vyhodnocuje reporty, riziká a plnenie termínov,
\item administrátor riadi prístupy, bezpečnosť a konfiguráciu systému.
\end{itemize}

Hlavným prínosom konceptu je prepojenie prevádzkových procesov do jedného toku údajov: od evidencie zariadenia, cez plánovanú revíziu, až po servisný zásah a reportovanie pre audit.

\section{Modulový koncept riešenia}

Systém je rozdelený do modulov tak, aby bolo možné riešenie iteratívne implementovať a rozširovať:

\begin{enumerate}[label=\arabic*.]
\item \textbf{Register zariadení} -- evidencia zariadení, ich parametrov, väzieb a histórie zmien.
\item \textbf{Revízie a harmonogram} -- plánovanie revízií, evidencia výsledkov, protokolov a termínov.
\item \textbf{Servis a incidenty} -- založenie zásahu, priebeh riešenia, dokumentácia opráv.
\item \textbf{Notifikácie} -- upozornenia pred termínom a po termíne, konfigurácia príjemcov.
\item \textbf{Reporty a export} -- prevádzkové a manažérske výstupy, export do PDF/CSV.
\item \textbf{Bezpečnosť a audit} -- autentifikácia, autorizácia, auditný log kritických operácií.
\item \textbf{Integrácie} -- príprava rozhraní pre prepojenie s GIS, SCADA a CMMS/EAM.
\end{enumerate}

Takto definované moduly umožňujú postupné dodávanie funkcií bez potreby zásadnej zmeny architektúry.

\section{Odhad veľkosti systému}

Pre potreby prvej iterácie bol vytvorený hrubý odhad veľkosti na úrovni modulov. Odhad vyjadruje relatívnu komplexnosť a predpokladaný rozsah implementácie (MVP + základ pre rozšírenie).

\begin{table}[htbp]
\centering
\caption{Hrubý odhad veľkosti podľa modulov}
\label{tab:rough-size-estimate}
\small
\begin{tabular}{p{4.2cm} c c p{4.3cm}}
\toprule
Modul & Komplexnosť & Odhad SP & Poznámka k rozsahu \\
\midrule
Register zariadení & stredná & 18 & CRUD, história zmien, filtrovanie \\
Revízie a harmonogram & vysoká & 24 & plánovanie, stavy, prílohy, termíny \\
Servis a incidenty & stredná až vysoká & 20 & workflow zásahov, väzby na revízie \\
Notifikácie & stredná & 12 & pravidlá, e-mail, interné upozornenia \\
Reporty a export & stredná & 14 & prehľady, PDF/CSV export \\
Bezpečnosť a audit & vysoká & 18 & roly, oprávnenia, auditné záznamy \\
Integrácie (API vrstva) & vysoká & 22 & príprava konektorov, mapovanie dát \\
\midrule
\textbf{Spolu} &  & \textbf{128 SP} & hrubý odhad pre jadro systému \\
\bottomrule
\end{tabular}
\end{table}

Na základe tabuľky je jadro systému odhadované na približne 128 story points. Pri stabilnej kapacite tímu 28--36 SP na iteráciu to predstavuje 4 až 5 iterácií pre dodanie funkčného jadra vrátane testovania a základnej dokumentácie.

Orientačný objem implementácie pre prvú plnohodnotnú verziu jadra systému je:

\begin{itemize}
\item \textbf{backend + integrácie}: približne 9\,000 až 13\,000 riadkov kódu,
\item \textbf{frontend}: približne 4\,000 až 6\,000 riadkov kódu,
\item \textbf{testy a skripty}: približne 2\,000 až 3\,500 riadkov kódu.
\end{itemize}

Celkový odhad sa pohybuje v intervale približne 15\,000 až 22\,500 riadkov kódu. Ide o plánovací odhad, ktorý bude spresňovaný po detailnom návrhu databázového modelu, API kontraktov a po implementácii prvých sprintov.

\subsection{Odhad metódou FPA a COCOMO}

Popri modulovom odhade bol pre validáciu veľkosti použitý aj orientačný odhad cez Function Point Analysis (FPA) a model COCOMO.

\subsubsection*{Výpočet UFP}

Tabuľka~\ref{tab:ufp} uvádza rozčlenenie funkčných bodov podľa typu a komplexnosti. Váhy zodpovedajú štandardu IFPUG pre semi-detached systémy (EI: 3/4/6, EO: 4/5/7, EQ: 3/4/6, ILF: 7/10/15, EIF: 5/7/10).

\begin{table}[htbp]
\centering
\caption{Odhad Unadjusted Function Points (UFP)}
\label{tab:ufp}
\small
\begin{tabular}{l c c c r}
\toprule
Typ funkčného bodu & Jednoduchá & Priemerná & Zložitá & UFP \\
\midrule
Externé vstupy (EI) & 3 & 3 & 2 & 33 \\
Externé výstupy (EO) & 2 & 2 & 0 & 18 \\
Externé dopyty (EQ) & 2 & 2 & 1 & 20 \\
Interné logické súbory (ILF) & 1 & 2 & 0 & 27 \\
Externé rozhrania (EIF) & 3 & 0 & 0 & 15 \\
\midrule
\textbf{Spolu (UFP)} & & & & \textbf{113} \\
\bottomrule
\end{tabular}
\end{table}

\subsubsection*{Výpočet VAF}

Value Adjustment Factor (VAF) sa vypočíta z celkového stupňa vplyvu (TDI) 14 všeobecných charakteristík systému (GSC), kde každá je hodnotená na škále 0--5. Tabuľka~\ref{tab:vaf} uvádza zvolené hodnoty pre navrhovaný systém.

\begin{table}[htbp]
\centering
\caption{Hodnotenie všeobecných charakteristík systému (GSC) a výpočet VAF}
\label{tab:vaf}
\small
\begin{tabular}{c l c}
\toprule
\# & Charakteristika (GSC) & DI (0--5) \\
\midrule
1 & Dátová komunikácia & 4 \\
2 & Distribuované spracovanie dát & 3 \\
3 & Výkon & 4 \\
4 & Intenzívne využívanie konfigurácie & 3 \\
5 & Frekvencia transakcií & 3 \\
6 & Online zadávanie dát & 4 \\
7 & Efektivita pre koncových používateľov & 4 \\
8 & Online aktualizácia & 4 \\
9 & Komplexné spracovanie & 3 \\
10 & Opätovná využiteľnosť & 2 \\
11 & Jednoduchosť inštalácie & 2 \\
12 & Jednoduchosť prevádzky & 3 \\
13 & Viacero lokalít nasadenia & 2 \\
14 & Podpora zmien (rozšíriteľnosť) & 5 \\
\midrule
 & \textbf{TDI ($\Sigma$ DI)} & \textbf{46} \\
 & \textbf{VAF $= 0{,}65 + 0{,}01 \times 46$} & \textbf{1{,}11} \\
\bottomrule
\end{tabular}
\end{table}

Z toho vyplýva: $\text{AFP} = \text{UFP} \times \text{VAF} = 113 \times 1{,}11 \approx 125$ FP. Pri priemernej hustote kódu $\sim$70 LOC/FP (webová aplikácia) to zodpovedá orientačne $\sim$8\,750 LOC, čo je vstupný objem pre výpočet COCOMO.

Pre Intermediate COCOMO bol dodatočne uplatnený faktor úpravy námahy $\text{EAF} = 1{,}22$, ktorý zohľadňuje zvýšenú zložitosť bezpečnostných požiadaviek a integrácií. Základný (Basic) COCOMO a Intermediárny (Intermediate) COCOMO poskytli tento orientačný výsledok:

\begin{table}[htbp]
\centering
\caption{Zhrnutie odhadu Basic COCOMO a Intermediate COCOMO}
\label{tab:cocomo-summary}
\small
\begin{tabular}{p{4.3cm} c c c c}
\toprule
Model & E & D & P & EAF \\
\midrule
Basic COCOMO & 34,1 & 8,6 & 4 & -- \\
Intermediate COCOMO & 41,6 & 9,2 & 5 & 1,22 \\
\bottomrule
\end{tabular}
\end{table}

Kde $E$ predstavuje odhad námahy (person-month), $D$ trvanie projektu (mesiace), $P$ veľkosť tímu a $EAF$ faktor úpravy námahy. Tento odhad je určený na plánovanie iterácií a kapacít tímu, nie ako pevný záväzný harmonogram.

\subsection{Person-months, veľkosť tímu a náklady}

Nasledujúci výpočet sumarizuje konkrétne hodnoty person-months (PM), odporúčanú veľkosť tímu pri plánovanom trvaní projektu a orientačné náklady.

\begin{itemize}
\item Z nameraných COCOMO hodnôt: \textbf{Basic COCOMO} $E_{basic}=34{,}1$ PM, \textbf{Intermediate COCOMO} $E_{int}=41{,}6$ PM.
\item Predpokladané trvanie projektu: $D=9$ mesiacov (zaokrúhlené z $D_{basic}=8{,}6$ a $D_{int}=9{,}2$ vypočítaných COCOMO).
\end{itemize}

Veľkosť tímu (hrubý odhad) počítame ako $\mathrm{team}=\lceil E / D \rceil$ (zaokrúhlené nahor):

\begin{itemize}
\item Basic COCOMO: $\mathrm{team}_{basic}=\lceil 34{,}1 / 9 \rceil = 4$ členovia,
\item Intermediate COCOMO: $\mathrm{team}_{int}=\lceil 41{,}6 / 9 \rceil = 5$ členovia.
\end{itemize}

Orientačné náklady možno odhadnúť pri predpokladanej trhovej sadzbe 3\,000\,–\,4\,000\,€ za PM (junior--mid vývojár vrátane réžie):

\begin{itemize}
\item Basic COCOMO ($E_{basic}=34{,}1$ PM): náklady $\approx 102\,000\,–\,136\,000$ €,
\item Intermediate COCOMO ($E_{int}=41{,}6$ PM): náklady $\approx 125\,000\,–\,166\,000$ €.
\end{itemize}

Poznámky a odporúčania:
\begin{itemize}
\item Pre konzervatívne plánovanie odporúčame použiť Intermediate COCOMO ($E_{int}=41{,}6$ PM) z dôvodu rizík integrácií a bezpečnosti.
\item Navrhovaná štruktúra tímu pre $4$--$6$ členov: 1 projektový/produktový manažér (časť úväzku), 1--2 backend developeri, 1 frontend developer, 1 QA/tester, 1 integračný/DevOps špecialista (pre 6-členný tím možno kombinovať role podľa potreby).
\item Uvádzané náklady sú orientačné a nezahŕňajú režijné náklady (infra, licencie, externé služby). Pri detailnom rozpočte odporúčame zohľadniť 10--20\% režijnú rezervu.
\end{itemize}

